\documentclass{article}
\usepackage[utf8]{inputenc}
\usepackage[spanish]{babel}
\usepackage{amsmath, amsthm, amssymb}
\usepackage[margin=1in]{geometry}

\renewcommand{\qedsymbol}{$\blacksquare$}

\begin{document}

\begin{enumerate}
    \item \textbf{Recurrencia $T(n) = 4T(n/2) + n$}
    
    Utilice el método de substitución para determinar la solución a la siguiente recurrencia. La solución de acuerdo con el \textit{Master Method} es $\Theta(n^2)$, pero usar la hipótesis $cn^2$ falla.
    
    \paragraph{1.1 Intento con hipótesis $T(n) \le cn^2$}
    Asumimos que $T(k) \le ck^2$ para todo $k < n$. Sustituyendo en la recurrencia:
    \begin{align*}
        T(n) &= 4T(n/2) + n \\
        T(n) &\le 4c(n/2)^2 + n \\
        T(n) &\le 4c(n^2/4) + n \\
        T(n) &\le cn^2 + n
    \end{align*}
    Para que la inducción sea válida, se requiere que $cn^2 + n \le cn^2$, lo cual implica $n \le 0$. Como estamos analizando el crecimiento para $n \ge 1$, la hipótesis \textbf{falla}.

    \paragraph{1.2 Hipótesis modificada: $T(n) \le cn^2 - dn$}
    Para compensar el término lineal sobrante, asumimos $T(k) \le ck^2 - dk$:
    \begin{align*}
        T(n) &\le 4(c(n/2)^2 - d(n/2)) + n \\
        T(n) &= 4(cn^2/4 - dn/2) + n \\
        T(n) &= cn^2 - 2dn + n
    \end{align*}
    Queremos demostrar que $cn^2 - 2dn + n \le cn^2 - dn$. Simplificando la desigualdad:
    \begin{align*}
        -2dn + n &\le -dn \\
        n &\le dn \\
        d &\ge 1
    \end{align*}
    La desigualdad se cumple para cualquier $d \ge 1$. Por lo tanto, se confirma que $T(n) = \Theta(n^2)$. \qed

    \vspace{2em}

    \item \textbf{Recurrencia $T(n) = 3T(\sqrt{n}) + \log_2 n$}
    
    \paragraph{2.1 Cambio de variable}
    Sea $n = 2^m$, lo que implica que $m = \log_2 n$ y $\sqrt{n} = 2^{m/2}$. Sustituimos en la recurrencia original:
    \[ T(2^m) = 3T(2^{m/2}) + m \]
    Definimos una nueva función $S(m) = T(2^m)$, obteniendo la forma:
    \[ S(m) = 3S(m/2) + m \]

    \paragraph{2.2 Método de substitución para $S(m)$}
    Proponemos como hipótesis (dada la forma de la relación de recurrencia) $S(k) \le ck^{\log_2 3} - dk$ para compensar el término lineal $m$:
    \begin{align*}
        S(m) &\le 3(c(m/2)^{\log_2 3} - d(m/2)) + m \\
        S(m) &= 3\left(c \frac{m^{\log_2 3}}{3} - \frac{dm}{2}\right) + m \\
        S(m) &= cm^{\log_2 3} - \frac{3}{2}dm + m
    \end{align*}
    Buscamos que $cm^{\log_2 3} - \frac{3}{2}dm + m \le cm^{\log_2 3} - dm$:
    \begin{align*}
        -\frac{3}{2}dm + m &\le -dm \\
        m &\le \frac{1}{2}dm \\
        d &\ge 2
    \end{align*}
    Esto demuestra que $S(m) = O(m^{\log_2 3})$.

    \paragraph{2.3 Regreso a la variable original}
    Dado que $S(m) = T(n)$ y $m = \log_2 n$:
    \[ T(n) = \Theta((\log_2 n)^{\log_2 3}) \approx \Theta((\log_2 n)^{1.58}) \] \qed
\end{enumerate}

\end{document}
